\subsection*{Matched supervision of visual question answering}
We joined pictured human questions and answers to their original training documents and selected a window of at most 192 characters by question overlap, using no reference answer to choose the window. Of 192,230 records, 715 windows overlapped held-out content, 245 contained unrenderable characters and 263 had no matching window. The 191,007 retained windows comprised 15,718,278 pictured positions. Retrieval can omit the evidence a correct answer needs. We ran one pass through this corpus and six passes through the original 2,565 examples for each arm. Fixed random derangements reassigned questions within each training corpus. The same passages remained available to both arms.

A shared visual encoder applied learned projections to pooled pixels. Dilated convolutions represented local question and passage context. A recurrent decoder with attention over both images emitted one bitmap or stopped. For a target bitmap $g$, let $a_w,a_c$ denote writing and copying probabilities, $w_g$ its probability among training-answer shapes and $c_j$ the conditional pointer probability of input position $j$ with bitmap $x_j$. We maximized the log probability of the physical outcome, $a_w w_g+a_c\sum_{j:x_j=g}c_j$, and the final stop probability. Inference made greedy decisions under this same distribution using only previous emissions. Training used 2,152 updates per arm and took 40.4 and 40.6 seconds on one RTX 4090 D, excluding preparation, evaluation and control runs. Peak PyTorch allocations were 1.04\,GiB per fit. All 14,860 development targets, including stops, were reachable with reference histories. We directly recounted the outputs of all four inference conditions. Exact reference agreement can undercount valid paraphrases.

For the subsequent recognition test, we paired held-out examples within each answer group by increasing reference length, shuffling ties with a fixed seed. One case lacked a partner. Each trial supplied the two passages in a common strip, one question and both candidate responses. We scored both responses under their own complete histories using mean log probability per target, including the stop. Reversing the question and the passage order gave four decisions per pair. The lexical controls used the Jaccard similarity between two sets of physical bitmaps or the maximum product of question-to-passage and passage-to-answer similarities over both regions. Ties received no credit.
